\begin{textsample}{NEG Dim 1 – persona\_gemini – Score -1.00 – t081\_gemini.txt}  \label{ex:f1_neg_020}
Indigenous Peoples comprise less than 5 % of the global population, yet they are responsible for protecting 80 % of the world 's biodiversity.
This is achieved despite ongoing rights \textbf{violations} and their exclusion from critical decisions.
At COP15, Indigenous leaders are speaking out.
Valentin Engobo from the Democratic Republic of Congo is highlighting the success of community forestry and recent peatland discoveries that store vast amounts of carbon.
Orpha from West Papua is fighting oil \textbf{palm} \textbf{plantations} that are clearing forests without community consent.
In Canada, Adrienne Jerome is advocating for the protection of caribou, which are under threat from industrial activities.
From Brazil, Dinamam Tuxá is denouncing policies that dispossess Indigenous Peoples of their territories.
The message to COP15 is clear : for biodiversity protection to be effective, colonial conservation models must be abandoned.
It is essential to ensure the full participation of Indigenous Peoples and to recognize their rights.

% matched lemmas: palm, plantation, violation
\end{textsample}
